\subsection{Sensitivity to serving batch size}
\label{sec:batch_sensitivity}

The serving results above are reported at a decode batch of $32$. Batch is a
workload parameter rather than a decision variable, but it is not a neutral one:
it reshapes the program's coefficients in four places at once, namely the
mixture-of-experts reload traffic through the number of distinct experts a step
touches, the decode compute, the KV read, and the array utilization of the
decode operator. We therefore rebuild the program at
$B \in \{1, 8, 32, 128\}$ and re-certify, comparing in each case against a
TPU-v3 baseline evaluated at the \emph{same} batch, so that the comparison is
like for like at every operating point.

\Cref{tab:batch} reports the result. Two things follow, and the second is
specific to mixture-of-experts serving.

\paragraph{The advantage is not an artifact of the operating point, and
$B=32$ is conservative.}
The certified serving speedup is $33.1\times$ at $B=1$ and rises monotonically
to $48.6\times$ at $B=128$. The claim of a large serving advantage therefore
does not depend on a favorable choice of batch: it holds from single-stream
interactive decoding to throughput-oriented batch serving, and the batch we
report elsewhere sits below the most favorable point rather than at it. The
increase with batch has a simple cause: the fixed baseline must stream a reload
that grows with batch across a bandwidth it cannot change, whereas the
co-designed chip provisions bandwidth and applies the byte-reducing decisions, so
the gap widens exactly where the baseline is most starved.

\paragraph{Batching is a weak amortizer for a sparse mixture of experts, which
is why speculation retains its value.}
For a dense model, a larger batch amortizes each weight load perfectly: the same
weights serve every token in the batch. Top-$k$ routing breaks this. The
expected number of distinct experts touched per step follows
$n(1 - (1 - k/n)^B)$, so a larger batch touches \emph{more} experts and the
reload grows: from $19$~GB per step at $B=1$ to $598$~GB at $B=128$ in
\cref{tab:batch}. Per token the reload does fall, from $19$ to $4.7$~GB, but
this is a $4\times$ amortization for a $128\times$ increase in batch, and it
saturates once essentially every expert is swept ($251.6$ of $256$ at $B=128$).
The consequence is that the decode memory wall is not a low-batch phenomenon
that disappears under load, and the optimal speculative width does not
collapse as batch grows but \emph{increases} with it: $K^\star$ doubles with
every fourfold increase in batch, from $2$ at $B=1$ through $8$ and $16$ to
$32$ at $B=128$. The reason is visible in the same table: because the reload
grows with batch, each verification pass has more traffic to amortize, so wider
speculation keeps paying. Speculation and batching are commonly described as
competing amortizers; for a sparse mixture of experts they are better described
as complementary, since batching amortizes the reload far less effectively than
it does for a dense model and thereby leaves the work that speculation does.

\begin{table}[h]
\centering
\footnotesize
\caption{Certified serving speedup against a TPU-v3 baseline evaluated at the
same batch ($p = 0.75$, $1$~GB diffusion draft). $\mathbb{E}[\text{experts}]$ is
the expected number of distinct routed experts touched per decode step, out of
$256$. Best-of-$K$ is the certified achievable design at that batch. All
entries, including the baselines, come from one search setting
(\cref{app:coeff_provenance}): a $24$-point grid on the compute reciprocal, an
integral bandwidth branch, and full enumeration of the precision, sparsity,
memory-technology and array-aspect selectors.}
\label{tab:batch}
\begin{tabular}{@{}rrrrrrrl@{}}
\toprule
$B$ & $\mathbb{E}[\text{experts}]$ & Reload (GB/step) & $K{=}2$ & $K{=}8$ &
$K{=}16$ & $K{=}32$ & Best \\
\midrule
$1$   & $8.0$   & $19.0$  & $\mathbf{33.13}\times$ & $30.12\times$ & $23.28\times$ & $15.69\times$ & $33.13\times$ at $K^\star{=}2$ \\
$8$   & $57.4$  & $136.6$ & $28.35\times$ & $\mathbf{36.76}\times$ & $32.18\times$ & $24.32\times$ & $36.76\times$ at $K^\star{=}8$ \\
$32$  & $163.3$ & $388.5$ & $28.56\times$ & $43.03\times$ & $\mathbf{44.85}\times$ & $36.01\times$ & $44.85\times$ at $K^\star{=}16$ \\
$128$ & $251.6$ & $598.5$ & $29.02\times$ & $46.11\times$ & $48.39\times$ & $\mathbf{48.64}\times$ & $48.64\times$ at $K^\star{=}32$ \\
\bottomrule
\end{tabular}
\end{table}

\subsection{Performance per TDP}
\label{sec:perf_tdp}

Speedup alone does not settle whether a design is worth building: a chip that is
fast because it is large is not obviously an advance. FAST~\citep{zhang_fast}
therefore searches on performance per unit power, and reports a headline of
about $3.7\times$ over its baseline. We evaluate the same quantity on the
designs of \cref{tab:batch}, using one power model for both sides
(\cref{app:coeff_provenance}) so that the ratio is convention-consistent:
$\mathrm{Perf/TDP} = \text{speedup} / (\mathrm{TDP}_{\text{design}} /
\mathrm{TDP}_{\text{baseline}})$, with the TPU-v3 baseline at $41.9$~W.

\Cref{tab:perftdp} reports the result. 
Performance per TDP rises from
$2.73\times$ at $B=1$ to $4.83\times$ at $B=128$, exceeding the FAST headline
for batches of $32$ and above and approaching it at $B=8$. Two points are worth
making precisely. First, the designs are selected by \emph{serving latency},
not by performance per TDP, so this is the efficiency of a latency-optimal
design rather than the best efficiency the program can reach; searching the same
certified frontier on performance per TDP directly moves the figure by less than
$2\%$ at $B \le 8$ and selects the identical design at $B=32$, which indicates
that the speedup is not being bought with disproportionate power. Second, the
efficiency improves with batch for the same structural reason the speedup does:
the fixed baseline must stream a reload that grows with batch across bandwidth
it cannot change, while the co-designed chip holds a single $24.4\times$-MAC,
$3200$~GB/s configuration across $B \in \{8, 32, 128\}$ and simply serves more
tokens per unit of that fixed power.

\begin{table}[h]
\centering
\footnotesize
\caption{Performance per TDP of the certified designs of \cref{tab:batch}, in
the FAST convention (speedup normalized by the TDP ratio; TPU-v3 baseline
$41.9$~W). FAST reports approximately $3.7\times$.
}
\label{tab:perftdp}
\begin{tabular}{@{}rrrrrrr@{}}
\toprule
$B$ & $K^\star$ & Speedup & MACs & BW (GB/s) & TDP (W) & Perf/TDP \\
\midrule
$1$   & $2$  & $33.13\times$ & $31.1\times$ & $3200$ & $509$ & $2.73\times$ \\
$8$   & $8$  & $36.76\times$ & $24.4\times$ & $3200$ & $422$ & $3.65\times$ \\
$32$  & $16$ & $44.85\times$ & $24.4\times$ & $3200$ & $422$ & $\mathbf{4.45}\times$ \\
$128$ & $32$ & $48.64\times$ & $24.4\times$ & $3200$ & $422$ & $\mathbf{4.83}\times$ \\
\bottomrule
\end{tabular}
\end{table}

\subsection{Certified bounds and sensitivity to draft acceptance}
\label{sec:certified_gap}

Every design reported above is a certified achievable point: a concrete,
feasibility-checked configuration whose serving time we can exhibit. Solving the
same program with the physically integral selectors relaxed gives the
complementary quantity, a bound that no design in this model can beat. Together
they bracket the unknown optimum, and the width of the bracket is a certified
statement about how much performance could still be left on the table. This is
the quantity a point estimate cannot provide, and we report it under one search
configuration so that both ends come from the same search.

\Cref{tab:gap} gives the bracket across batch and draft acceptance. Three
readings follow. First, \emph{the bracket tightens where throughput serving
lives}: the gap falls from $63.8\%$ at $B=1$ to $13.5\%$ at $B=128$ with
$p=0.60$, so in the high-batch regime the optimum is pinned to within about a
seventh of the achievable value. Second, \emph{acceptance matters unevenly}. At
$B=1$ the three acceptance columns are identical, because $K^\star=2$ consumes
almost no acceptance; at $B=128$ the achievable spans $34.6\times$ to
$66.5\times$ as $p$ moves from $0.60$ to $0.85$. The parameter we state rather
than measure is therefore nearly inert in the latency regime and dominant in the
throughput regime, which is where a measured acceptance rate would sharpen the
result most. Third, \emph{the relaxation loosens as acceptance rises} (from
$+40.0\%$ to $+59.0\%$ at $B=8$): a higher acceptance rate gives the fractional
relaxation more room to exploit, so the bound degrades in exactly the regime
where the achievable improves, and we report both rather than only the
favourable half.

We note that the operating point used elsewhere in this paper, $p=0.75$, is the
middle of the swept range: the $p=0.85$ column is materially higher at every
batch, so the headline figures are conservative on the acceptance axis as well
as on the batch axis.

\begin{table}[h]
\centering
\footnotesize
\caption{Certified bracket on the serving optimum under one search configuration.
``Achv'' is the certified achievable design, ``bound'' the relaxed-LP bound that
no design in this model can beat, and ``gap'' their separation. Both ends are
computed by the same search, and the $B=32$, $p=0.75$ entry reproduces
\cref{tab:batch} exactly.}
\label{tab:gap}
\begin{tabular}{@{}rr rrr rrr rrr@{}}
\toprule
& & \multicolumn{3}{c}{$p=0.60$} & \multicolumn{3}{c}{$p=0.75$} & \multicolumn{3}{c}{$p=0.85$} \\
\cmidrule(lr){3-5}\cmidrule(lr){6-8}\cmidrule(lr){9-11}
$B$ & $K^\star$ & Achv & Bound & Gap & Achv & Bound & Gap & Achv & Bound & Gap \\
\midrule
$1$   & $2$  & $32.93\times$ & $53.96\times$ & $63.8\%$ & $33.13\times$ & $54.27\times$ & $63.8\%$ & $33.13\times$ & $54.27\times$ & $63.8\%$ \\
$8$   & $8$  & $31.32\times$ & $43.86\times$ & $40.0\%$ & $36.76\times$ & $55.28\times$ & $50.4\%$ & $40.72\times$ & $64.74\times$ & $59.0\%$ \\
$32$  & $16$ & $33.46\times$ & $39.91\times$ & $19.3\%$ & $44.85\times$ & $57.26\times$ & $27.7\%$ & $56.85\times$ & $78.38\times$ & $37.9\%$ \\
$128$ & $32$ & $34.60\times$ & $39.26\times$ & $\mathbf{13.5\%}$ & $48.64\times$ & $58.37\times$ & $20.0\%$ & $66.49\times$ & $86.13\times$ & $29.5\%$ \\
\bottomrule
\end{tabular}
\end{table}

\subsection{Does sparse routing cause the effect?}
\label{sec:dense_control}

The explanation offered above is causal: we claim that the widening of
$K^\star$ and the growth of the advantage with batch are consequences of top-$k$
routing, which makes the reload grow with batch instead of amortizing. That
claim is testable by ablation. We rebuild the identical program with one line
changed, pinning the expected distinct-expert count to its $B=1$ value so that a
larger batch reuses the same expert set and weight loads amortize perfectly, as
they do for a dense model. Model dimensions, KV traffic, decode compute, the
decision menus, the speculative-width menu, the search setting and the baseline
convention are untouched, so any difference isolates routing and nothing else.
The prediction is sharp: under dense-equivalent routing both effects should
disappear.

\Cref{tab:dense} reports the comparison. They disappear. With the reload pinned
at $19.0$~GB per step, the selected speculative width is $K^\star = 2$ at every
batch and the certified speedup is flat near $33\times$, against
$K^\star = 2, 8, 16, 32$ and $33.1\times$ to $48.6\times$ under real routing.
The ablation also supplies its own positive control: at $B=1$ sparse and
dense routing coincide by construction, and the two arms agree to four
significant figures ($33.13\times$, $K^\star=2$ in both), which confirms that
the control differs from the treatment in the intended variable only.

We therefore state the mechanism as measured rather than argued. Speculation
scales with batch in sparse serving because each verification pass has a larger
reload to amortize; remove that growth and the optimal width collapses to its
low-batch value. The practical reading is that a serving stack tuned on dense
models will systematically under-speculate on a mixture of experts, and that the
error grows with the batch size it was tuned at.

\begin{table}[h]
\centering
\footnotesize
\caption{Controlled ablation of sparse routing. The dense-control column pins the
distinct-expert count to its $B=1$ value so weight loads amortize perfectly with
batch; everything else is identical, including the search setting. Both
signature effects vanish, and the two arms agree exactly at $B=1$, where the
ablation is a no-op.}
\label{tab:dense}
\begin{tabular}{@{}rrrrrr@{}}
\toprule
& \multicolumn{2}{c}{Sparse (top-$k$ routing)} & & \multicolumn{2}{c}{Dense control} \\
\cmidrule(lr){2-3}\cmidrule(lr){5-6}
$B$ & Reload (GB) & $K^\star$ (speedup) & & Reload (GB) & $K^\star$ (speedup) \\
\midrule
$1$   & $19.0$  & $2$  ($33.13\times$) & & $19.0$ & $2$ ($33.13\times$) \\
$8$   & $136.6$ & $8$  ($36.76\times$) & & $19.0$ & $2$ ($33.08\times$) \\
$32$  & $388.5$ & $16$ ($44.85\times$) & & $19.0$ & $2$ ($33.06\times$) \\
$128$ & $598.5$ & $32$ ($48.64\times$) & & $19.0$ & $2$ ($33.24\times$) \\
\bottomrule
\end{tabular}
\end{table}

